\documentclass[../../main.tex]{subfiles}

\begin{document}

\section{Stable Arithmetic Equational Spectra and Affine-Linear Audit}\label{app:stable-arithmetic-equational-audit}

This appendix records the certificate boundary for the stable arithmetic
operations when they are viewed as finite magmas in the Equational Theories
Project database.  Let
\[
    \mathcal E=\{E_1,\ldots,E_{4694}\}
\]
denote the current order-four ETP equation list at commit
\[
    \texttt{b99c0e486501e5b0690ef3fe5250d3aa57e63478}.
\]
The positive result is a clean affine-linear transport theorem and a complete
computed spectrum audit.  The negative result is equally important: neither
the stable arithmetic spectra nor the full Fibonacci affine-linear family meet
any of the current dashboard unknowns.  Thus this appendix is an audit and
reproducibility record, not a claim of a new ETP dashboard resolution.

\subsection{Transport and coefficient criteria}

For a magma term \(t\) in the variables
\[
    x,y,z,w,u,v,
\]
write \(\nu(t)\in\NN^6\) for its variable multiplicity vector.  If a binary
operation is affine-linear over \(\ZZ/n\ZZ\),
\[
    x\diamond y=ax+by,
\]
then every term \(t\) evaluates to a linear form
\[
    t=\sum_i \lambda_i(t;n,a,b)x_i
\]
with coefficients in \(\ZZ/n\ZZ\), obtained by the recursion
\[
    \lambda(u\diamond v;n,a,b)
       =a\,\lambda(u;n,a,b)+b\,\lambda(v;n,a,b).
\]

\begin{lemma}[Affine-linear coefficient criterion]\label{lem:stable-audit-affine-coefficient-criterion}
For the affine-linear magma \(x\diamond y=ax+by\) on \(\ZZ/n\ZZ\), an ETP
equation \(s=t\) holds for all assignments if and only if
\[
    \lambda(s;n,a,b)=\lambda(t;n,a,b)
        \quad\text{in }(\ZZ/n\ZZ)^6.
\]
\end{lemma}

\begin{proof}
The difference \(s-t\) is a linear form in the six independent input variables
over \(\ZZ/n\ZZ\).  If all coefficient differences vanish, the equation holds
for every assignment.  Conversely, if the \(i\)-th coefficient difference is
nonzero, evaluating at \(x_i=1\) and all other variables \(0\) separates the two
sides.
\end{proof}

\begin{lemma}[Stable-value affine transfer]\label{lem:stable-audit-stable-value-affine-transfer}
Let \(X_m\) be the stable fiber and let
\[
    X_m\simeq \ZZ/F_{m+2}\ZZ
\]
be the stable-value ring equivalence.  For every affine-linear operation
\[
    x\diamond y=ax+by
\]
on \(X_m\), satisfaction of any affine magma equation is equivalent to
satisfaction of the transported equation on \(\ZZ/F_{m+2}\ZZ\) with the same
coefficient formula.
\leanverified{stableValueRingEquiv\_preserves\_magma\_satisfaction}
\end{lemma}

\begin{proof}
The stable-value map is a ring equivalence and hence preserves addition,
multiplication, and scalar multiplication by stable coefficients.  Induction on
the local binary term syntax gives equality between evaluation in \(X_m\) and
evaluation after transport to \(\ZZ/F_{m+2}\ZZ\).  Injectivity of the
equivalence gives the equivalence of equation satisfaction.  The statement is
formalized as
\[
    \texttt{\detokenize{lean4/Omega/Folding/FiberRing.lean:216}}.
\]
\end{proof}

The two stable arithmetic operations over a prime field are the special cases
\[
    \mathrm{stableAdd}(x,y)=x+y,\qquad
    \mathrm{stableMul}(x,y)=xy.
\]
For \(\mathrm{stableAdd}\), the law \(s=t\) holds over \(\FF_p\) exactly when
\[
    \nu(s)\equiv \nu(t)\pmod p.
\]
For \(\mathrm{stableMul}\), the term \(t\) evaluates to the monomial
\[
    \prod_i x_i^{\nu_i(t)}.
\]
Thus \(s=t\) holds over \(\FF_p\) if and only if the two support sets agree and
all positive exponents are congruent modulo \(p-1\).  The criteria above were
verified computationally against exhaustive assignment evaluation for
\(\operatorname{Fin}(2)\) and \(\operatorname{Fin}(3)\) (see
\texttt{\detokenize{scripts/equational_theory/}}).

\subsection{Stable arithmetic spectra}

\begin{theorem}[Stable arithmetic ETP spectra]\label{thm:stable-audit-stable-arithmetic-spectra}
For the prime Fibonacci moduli \(2,3,5,13,89\), the ETP spectra of
\(\mathrm{stableAdd}\) and \(\mathrm{stableMul}\) are the spectra shown in
Table \ref{tab:stable-audit-stable-spectra}.  The generated union contains
\[
    3,157,020
\]
ordered anti-implications, of which
\[
    2,912,508
\]
are not covered by the current \(\texttt{full\_entries.json}\) finite
\(\texttt{Facts}\) entries.
\end{theorem}

\begin{table}[H]
\centering
\scriptsize
\setlength{\tabcolsep}{5pt}
\caption{Stable arithmetic spectra in the current ETP equation list.  The last
column counts ordered satisfied/refuted pairs not already covered by current
\(\texttt{full\_entries.json}\) \(\texttt{Facts}\) entries.}
\label{tab:stable-audit-stable-spectra}
\begin{tabular}{c r r r r r}
\toprule
operation & \(p\) & satisfied & refuted & anti-implications & not covered\\
\midrule
\(\mathrm{stableAdd}\) & 2  & 663 & 4031 & 2672553 & 2492999\\
\(\mathrm{stableMul}\) & 2  & 164 & 4530 & 742920  & 668079\\
\(\mathrm{stableAdd}\) & 3  & 60  & 4634 & 278040  & 249114\\
\(\mathrm{stableMul}\) & 3  & 93  & 4601 & 427893  & 386077\\
\(\mathrm{stableAdd}\) & 5  & 32  & 4662 & 149184  & 131647\\
\(\mathrm{stableMul}\) & 5  & 46  & 4648 & 213808  & 193915\\
\(\mathrm{stableAdd}\) & 13 & 32  & 4662 & 149184  & 131647\\
\(\mathrm{stableMul}\) & 13 & 32  & 4662 & 149184  & 131647\\
\(\mathrm{stableAdd}\) & 89 & 32  & 4662 & 149184  & 131647\\
\(\mathrm{stableMul}\) & 89 & 32  & 4662 & 149184  & 131647\\
\bottomrule
\end{tabular}
\end{table}

\begin{proof}
The addition rows apply the congruence criterion
\(\nu(s)\equiv\nu(t)\pmod p\).  The multiplication rows apply the support and
positive-exponent congruence criterion modulo \(p-1\).  Each row then gives a
truth vector on the \(4694\) equations; the number of ordered
anti-implications is the product of the satisfied and refuted counts for that
row, and the union counts are obtained by bitset union across all rows.  The
computation is verified computationally (see
\texttt{\detokenize{scripts/equational_theory/}}), specifically by
\[
    \texttt{\detokenize{scripts/equational_theory/stable_arithmetic_facts.py}}
\]
and the certificate
\[
    \texttt{\detokenize{scripts/equational_theory/stable_arithmetic_facts.json}}.
\]
The accompanying ETP Lean file
\[
    \texttt{\detokenize{scripts/equational_theory/stable_arithmetic_facts.lean}}
\]
contains the six non-duplicative \(\operatorname{Fin}(2)\),
\(\operatorname{Fin}(3)\), and \(\operatorname{Fin}(5)\) finite
\(\texttt{Facts}\) certificates and compiles with \(\texttt{decideFin!}\).
\end{proof}

\begin{corollary}[Dashboard disjointness of the stable spectra]\label{cor:stable-audit-stable-dashboard-disjoint}
The stable arithmetic anti-implication union has zero intersection with the
current dashboard unknowns.  In particular, the \(2,912,508\) anti-implications
not covered by current \(\texttt{full\_entries.json}\) resolve
\[
    0
\]
of the \(498\) current dashboard unknown ordered pairs.
\end{corollary}

\begin{proof}
The dashboard audit parses \(\texttt{home\_page/fme/unknowns.json}\), extracts
the \(498\) ordered pairs whose two sides are ETP equations, and intersects
those pairs with the generated stable arithmetic union and with the subunion
not covered by current finite \(\texttt{Facts}\) entries.  Both intersections
are empty.  This is verified computationally (see
\texttt{\detokenize{scripts/equational_theory/}}), by
\[
    \texttt{\detokenize{scripts/equational_theory/stable_dashboard_unknowns.py}}
\]
and its JSON certificate
\[
    \texttt{\detokenize{scripts/equational_theory/stable_dashboard_unknowns.json}}.
\]
\end{proof}

\subsection{Prime stabilization and the abelian-group shadow}

Let \(\mathcal B\subset\mathcal E\) be the set of balanced equations
\(\nu(s)=\nu(t)\).  In the current ETP equation list this set is
\[
\begin{aligned}
\mathcal B=\{&
1,43,4283,4290,4320,4358,4362,4364,4369,4380,4396,4398,\\
&4405,4433,4435,4442,4473,4480,4482,4512,4515,4525,\\
&4531,4541,4544,4598,4605,4635,4673,4677,4679,4684\}.
\end{aligned}
\]

\begin{theorem}[Prime stabilization with the \(p=5\) multiplication exception]\label{thm:stable-audit-prime-stabilization}
The operation \(\mathrm{stableAdd}\) satisfies exactly the \(32\) equations in
\(\mathcal B\) for every prime \(p\geq 5\).  The operation
\(\mathrm{stableMul}\) satisfies exactly \(\mathcal B\) for every prime
\(p\geq 7\), but not for \(p=5\).  Over \(\operatorname{Fin}(5)\),
\(\mathrm{stableMul}\) satisfies \(46\) equations, including the non-balanced
equation
\[
    E_{411}:\qquad x=x\diamond(x\diamond(x\diamond(x\diamond x))).
\]
\end{theorem}

\begin{proof}
For addition, satisfaction is the congruence
\(\nu(s)\equiv\nu(t)\pmod p\).  The largest absolute variable-multiplicity
difference occurring in the ETP database is \(5\).  Hence for \(p>5\) the
congruence forces equality.  The remaining prime \(p=5\) is checked directly:
no non-balanced difference vector has every coordinate divisible by \(5\).

For multiplication, satisfaction requires equal supports and equality of all
positive exponent differences modulo \(p-1\).  When the supports agree, the
largest absolute positive exponent difference in the ETP database is \(4\).
Thus \(p-1\geq 6\) forces exponent equality, and all primes \(p\geq 7\) give
exactly \(\mathcal B\).  At \(p=5\), exponent differences divisible by \(4\)
produce extra laws.  In particular \(E_{411}\) holds because \(x^5=x\) for all
\(x\in\FF_5\).  The computation is verified computationally (see
\texttt{\detokenize{scripts/equational_theory/}}), by
\[
    \texttt{\detokenize{scripts/equational_theory/stable_prime_stabilization.py}}
\]
and
\[
    \texttt{\detokenize{scripts/equational_theory/stable_prime_stabilization.json}}.
\]
\end{proof}

\begin{theorem}[Abelian-group identity classification]\label{thm:stable-audit-abelian-group-classification}
The set \(\mathcal B\) is exactly the set of ETP equations valid in every
abelian group under the interpretation
\[
    x\diamond y=x+y.
\]
Moreover, cyclic addition on \(\ZZ/4\ZZ\) satisfies \(116\) ETP equations,
namely the \(32\) equations in \(\mathcal B\) together with \(84\) torsion
identities, while cyclic addition on \(\ZZ/6\ZZ\) satisfies exactly the
\(32\) equations in \(\mathcal B\).
\end{theorem}

\begin{proof}
In an abelian group, every positive binary term evaluates to the sum of its
variables with natural-number multiplicities.  If \(\nu(s)=\nu(t)\), then
associativity and commutativity rewrite both terms to the same formal sum, so
the equation is valid in every abelian group.  If the multiplicity vectors
differ in the coordinate of a variable \(v\), the infinite cyclic group
\(\ZZ\) separates the equation by assigning \(v=1\) and all other variables
\(0\).  Therefore the universally valid equations are exactly the balanced
ones.  The extraction of the \(32\)-equation set and the finite cyclic checks
are verified computationally (see
\texttt{\detokenize{scripts/equational_theory/}}), by
\[
    \texttt{\detokenize{scripts/equational_theory/abelian_group_identity_check.py}}
\]
and
\[
    \texttt{\detokenize{scripts/equational_theory/abelian_group_identity_check.json}}.
\]
\end{proof}

\begin{remark}[Interpretation of the stable-addition stabilization]\label{rem:stable-audit-abelian-shadow}
The \(32\)-equation stabilization for \(\mathrm{stableAdd}\) is the standard
abelian-group identity criterion restricted to the finite ETP equation list,
together with the finite \(p=5\) torsion check in Theorem
\ref{thm:stable-audit-prime-stabilization}.  It is therefore a useful
consistency audit of the stable arithmetic interface, not a new universal
algebra theorem.
\end{remark}

\subsection{Full Fibonacci affine-linear dashboard closure}

\begin{theorem}[Fibonacci affine-linear dashboard closure]\label{thm:stable-audit-affine-fibonacci-dashboard-closure}
Let
\[
    N=\{2,3,5,8,13,21,34,55,89,144\}.
\]
For every \(n\in N\) and every pair \(a,b\in\ZZ/n\ZZ\), the affine-linear magma
\[
    x\diamond y=ax+by
\]
on \(\ZZ/n\ZZ\) resolves no current dashboard unknown.  Equivalently, after
checking all
\[
    \sum_{n\in N}n^2=33550
\]
parameter triples \((n,a,b)\), the intersection with the \(498\) current
dashboard unknown ordered pairs is empty.
\end{theorem}

\begin{table}[H]
\centering
\scriptsize
\setlength{\tabcolsep}{8pt}
\caption{Full Fibonacci affine-linear dashboard audit.  Each row checks all
\((a,b)\in(\ZZ/n\ZZ)^2\) against the current dashboard unknown pairs.}
\label{tab:stable-audit-affine-fibonacci-dashboard}
\begin{tabular}{r r r r}
\toprule
\(n\) & parameter pairs & spectra with dashboard hits & resolved dashboard pairs\\
\midrule
2   & 4     & 0 & 0\\
3   & 9     & 0 & 0\\
5   & 25    & 0 & 0\\
8   & 64    & 0 & 0\\
13  & 169   & 0 & 0\\
21  & 441   & 0 & 0\\
34  & 1156  & 0 & 0\\
55  & 3025  & 0 & 0\\
89  & 7921  & 0 & 0\\
144 & 20736 & 0 & 0\\
\bottomrule
\end{tabular}
\end{table}

\begin{proof}
By Lemma \ref{lem:stable-audit-affine-coefficient-criterion}, each equation is
decided by equality of two coefficient vectors modulo \(n\).  The scan applies
this criterion to the \(223\) distinct equation numbers appearing in the
\(498\) dashboard unknown ordered pairs and records whether a spectrum has an
antecedent true and consequent false for any unknown pair.  Table
\ref{tab:stable-audit-affine-fibonacci-dashboard} shows that no such spectrum
exists.  The coefficient-vector criterion was cross-checked by exhaustive
assignment evaluation for every parameter pair over \(\ZZ/2\ZZ\) and
\(\ZZ/3\ZZ\).  The computation is verified computationally (see
\texttt{\detokenize{scripts/equational_theory/}}), by
\[
    \texttt{\detokenize{scripts/equational_theory/affine_fibonacci_dashboard_scan.py}}
\]
and
\[
    \texttt{\detokenize{scripts/equational_theory/affine_fibonacci_dashboard_scan.json}}.
\]
\end{proof}

\subsection{Prime-field saturation redundancy}

\begin{proposition}[Redundancy of the older prime-field linear audit]\label{prop:stable-audit-prime-field-redundancy}
In the at-most-two-variable ETP sublist, the older prime-field linear scan
contains \(810\) equations and produces \(608948\) separated ordered pairs
after scanning primes through \(31\); the cumulative set is already saturated
by \(p=13\).  As anti-implication data, this scan is redundant with the
existing \texttt{\detokenize{_364}} prime-field baseline.
\end{proposition}

\begin{proof}
The recomputation uses the same coefficient-vector linearization as Lemma
\ref{lem:stable-audit-affine-coefficient-criterion}, restricted to ETP
equations involving at most two variables and to prime fields \(p\leq 31\).
The cumulative counts are
\[
\begin{array}{c|rrrrrrrrrrr}
p&2&3&5&7&11&13&17&19&23&29&31\\
\hline
\text{new}&439263&90422&56259&20008&2612&384&0&0&0&0&0\\
\text{cum.}&439263&529685&585944&605952&608564&608948&
608948&608948&608948&608948&608948.
\end{array}
\]
The \texttt{\detokenize{_364}} baseline uses all \(4694\) equations and prime fields through
\(199\), so every pair from the older at-most-two-variable, \(p\leq 31\) scan
is included by scope and prime-range inclusion.  The recomputation also checks
that the \(17588\) \(\ZZ/144\ZZ\) witnesses new after the
prime-fields-plus-\(\ZZ/8\ZZ\) baseline have zero overlap with the older scan.
This audit is verified computationally (see
\texttt{\detokenize{scripts/equational_theory/}}), by
\[
    \texttt{\detokenize{scripts/equational_theory/compare_prime_saturation_to_364.py}}
\]
and
\[
    \texttt{\detokenize{scripts/equational_theory/prime_saturation_vs_364.json}}.
\]
\end{proof}

\begin{remark}[Certificate boundary]\label{rem:stable-audit-certificate-boundary}
The Automath/Omega Lean contribution in this appendix is the stable-value
transport theorem of Lemma \ref{lem:stable-audit-stable-value-affine-transfer}.
The ETP indexed spectra, \(\texttt{Facts}\) lists, dashboard intersections,
prime-stabilization finite checks, and prime-field redundancy audit are finite
computational certificates.  They are included with this paper for
reproducibility under
\[
    \texttt{\detokenize{scripts/equational_theory/}}.
\]
\end{remark}

\end{document}
